\section{Input-dependent boundary instruments}\label{sec:cqboundary71}
Fix an orthonormal basis $e_1,\ldots,e_d$ of the input and the associated
complete dephasing map
$\mathcal D(X)=\sum_x\langle e_x,Xe_x\rangle|e_x\rangle\langle e_x|$.
For positive semidefinite $n\times n$ matrices $\sigma_{xy}$ satisfying
$\sum_y\operatorname{tr}\sigma_{xy}=1$ for every $x$, define
\begin{equation}\label{eq:cqfamily71}
 \Phi_{\boldsymbol\sigma,y}(X)
       =\sum_{x=1}^d\langle e_x,Xe_x\rangle\sigma_{xy},
       \qquad 1\le y\le m.
\end{equation}
These instruments measure a fixed input basis and then prepare an
outcome-dependent state. Their public outcome probabilities
$\sum_x\langle e_x,Xe_x\rangle\operatorname{tr}\sigma_{xy}$ may
depend on the input. Measurement destroys input coherences; the family
is not in general input erasing, and it is not the coherent tensor
family of Theorem~\ref{thm:coherententropy70}.

The Choi blocks and their ranks are
\begin{equation}\label{eq:cqchoi71}
 J_y=\bigoplus_{x=1}^d\sigma_{xy},\qquad
 \operatorname{rank}J_y=\sum_x\operatorname{rank}\sigma_{xy}.
\end{equation}
In particular the marginal constraint is exactly one unit-trace
condition per input row. Fix public integers $0\le r_{xy}\le n$,
with $\sum_y r_{xy}>0$ in every row, and let $\mathcal F_{\boldsymbol r}$
be the family with $\operatorname{rank}\sigma_{xy}\le r_{xy}$. Put
\begin{equation}\label{eq:cqdim71}
 v_x=\sum_y r_{xy}(2n-r_{xy})-1,\qquad V=\sum_x v_x.
\end{equation}
The ranks here are conditional output-state ranks, not independent
rank bounds on arbitrary Choi blocks. The input basis, dimensions and
rank bounds are fixed, not learned from an unknown device.

\begin{theorem}[Rank-boundary coding with input-dependent outcomes]\label{thm:cqentropy71}
Fix the preceding data and $0<\delta_*<2$. If $V>0$, the least covering
cardinality by arbitrary legal memoryless centres and the least
cardinality with rational decoded Choi blocks both obey
\begin{equation}\label{eq:cqentropy71}
 \log_2\mathcal M_N(\delta;\mathcal F_{\boldsymbol r})
   =\frac V2\log_2N+V\log_2(1/\delta)+O(1)
\end{equation}
uniformly for $N\ge1$ and $0<\delta\le\delta_*$. The error is the
unhalved final-state trace distance over arbitrary common quantum
testers with finite references, quantum memory, feedback and public
stopping bounded by $N$. A rational upper code has decoded centres in
\eqref{eq:cqfamily71}; it never increases any conditional block rank
and preserves every zero block exactly. Its encoder on rational target
data uses integer and rational arithmetic only. If $V=0$ the family
is a known singleton.

The remainder depends on the fixed dimensions, rank bounds and
$\delta_*$, but not on $N$, $\delta$ or a smallest positive output
eigenvalue. Payload length is distinct from encoder workspace and
expanded decoder storage. The decoded object is a mathematical
quantum instrument, not a physical classical simulator on an unknown
reference-entangled input.
\end{theorem}

\subsection{Centres and the adaptive metric}
Write $\Omega_x=\bigoplus_y\sigma_{xy}$ for the flagged state of one
input row. A useful general observation first identifies the allowed
centres.

\begin{proposition}[Idempotent input retraction]\label{prop:inputretraction71}
Let $\Pi$ be a fixed completely positive trace-preserving idempotent
input map and let every target satisfy $\Phi\Pi=\Phi$. Then
\begin{equation}\label{eq:inputretraction71}
             d_N(\Phi,K\Pi)\le d_N(\Phi,K)
\end{equation}
for every legal memoryless centre $K$ and every $N$. Thus covering
numbers with all legal centres equal those with centres fixed by
right composition with $\Pi$, at every radius.
For $\Pi=\mathcal D$, the retracted centre has rows
$\tau_{xy}=K_y(|e_x\rangle\langle e_x|)$ and is of the form
\eqref{eq:cqfamily71}. In the prescribed coordinate basis this operation
preserves rationality. It need not preserve the target's rank bounds.
\end{proposition}
\begin{proof}
For any tester of $\Phi$ against $K\Pi$, insert $\Pi$ on its input
immediately before each use of the unknown device. This is a legal
common tester of $\Phi$ against $K$: in the first hypothesis
$\Phi\Pi=\Phi$, and in the second it implements $K\Pi$. References,
feedback and stopping are unaffected. Taking the supremum proves
\eqref{eq:inputretraction71}. Idempotence says that $K\Pi$ is fixed
by right composition with $\Pi$. Map every covering centre this way;
cardinality cannot increase. The reverse covering inequality follows
from inclusion of allowed centre classes. For dephasing, the displayed
rows are positive and have unit total trace, since $K$ is an instrument.
They are diagonal input blocks of the original Choi matrices.
\end{proof}
This is a data-processing retraction, not a rank-preserving projection.
It generalizes the preparation-centre observation of
Theorem~\ref{thm:retraction70}, whose fixed input map discards its input
and supplies one chosen state. In the converse of
Theorem~\ref{thm:cqentropy71}, no rank promise is imposed on centres.

\begin{lemma}[Row witnesses and a common programme]\label{lem:cqmetric71}
For two instruments of \eqref{eq:cqfamily71},
\begin{equation}\label{eq:cqsandwich71}
 \max_x\|\Omega_x^{\otimes N}-\Lambda_x^{\otimes N}\|_1
 \le d_N(\Phi_{\boldsymbol\sigma},\Phi_{\boldsymbol\tau})
 \le\left\|\bigotimes_x\Omega_x^{\otimes N}
               -\bigotimes_x\Lambda_x^{\otimes N}\right\|_1.
\end{equation}
If unit purifications $a_x,b_x$ of the corresponding row states are
chosen and $\eta_x=\|a_x-b_x\|$, then
\begin{equation}\label{eq:cqfactorbound71}
 d_N(\Phi_{\boldsymbol\sigma},\Phi_{\boldsymbol\tau})
                  \le2\sqrt{N\sum_x\eta_x^2}.
\end{equation}
For one use there is the exact identity
$d_1=\max_x\|\Omega_x-\Lambda_x\|_1$.
\end{lemma}
\begin{proof}
For the first inequality repeatedly feed $|e_x\rangle\langle e_x|$
and retain all flagged outputs. For the second, prepare, for each
of the $N$ slots, one copy of every row state. A common processor
measures the actual input in the prescribed basis and uses only the
programme state with that outcome $x$. The remaining row states in
that slot are discarded. If the input is entangled with the tester's
memory, its conditional memory operator is retained by this measurement;
appending the selected row state gives exactly
\eqref{eq:cqfamily71} on that joint input. The unobserved input outcome
is internal to the processor, not new feedback granted to the tester.
All later operations and the discard of unused slots at a stop form one
common CPTP processor. Trace contraction proves the upper bound.

Purify all these programme states by the chosen vectors. For unit
vectors, $1-|\langle a_x,b_x\rangle|^2\le\eta_x^2$. The elementary
product inequality $1-\prod t_j\le\sum(1-t_j)$ for $t_j\in[0,1]$
gives pure-programme trace distance squared at most
$4N\sum_x\eta_x^2$. Partial trace and the common processor give
\eqref{eq:cqfactorbound71}. Finally, for an arbitrary one-use input
with reference, the difference is
$\sum_x\xi_R^x\otimes(\Omega_x-\Lambda_x)$, where
$\xi_R^x\ge0$ and $\sum_x\operatorname{tr}\xi_R^x=1$.
The triangle inequality gives the one-use upper maximum, and a basis
input attains it.
\end{proof}
The product programme consumes up to $dN$ row states in the proof,
whereas the actual tester makes at most $N$ instrument calls. It is
only a comparison device, not an equality of these resource counts.
Neither inequality in \eqref{eq:cqsandwich71} is asserted to be an
equality for general $N$. In particular the absence of an asymptotic
adaptive advantage in a binary discrimination exponent does not imply
a finite-use metric equality; compare \cite[Sections~II and V]{SHW68}.

\subsection{Exact rational decoding and the covering proof}
\begin{lemma}[A common-grid row code]\label{lem:cqcodec71}
For rational target data, apply the singular factor encoder of
Theorem~\ref{thm:factorcode68} separately in every input row at a
common integer grid $B\ge1$. If $h_x$ is that row's actual real
factor-coordinate count, the decoded instrument satisfies
\begin{equation}\label{eq:cqcodeerror71}
 d_N(\Phi,\widetilde\Phi)^2
       \le\frac{16N\sum_x(h_x-1)}{B^2}
       \le\frac{16NV}{B^2}.
\end{equation}
There are at most $C_{d,n,m,\boldsymbol r}(2B+1)^V$ legal decoder
images. The conditional block ranks do not increase; consequently no
outcome Choi rank increases either. All zero conditional blocks remain
zero. The expanded common denominator is at most
$B^{2d}\prod_x h_x$ and thus has $O_{d,n,m}(1+\log B)$ bits.
\end{lemma}
\begin{proof}
The inherited anchor construction gives a unit factor purification in
row $x$ within distance $2\sqrt{h_x-1}/B$ of the target's chosen
unit factor vector. Apply \eqref{eq:cqfactorbound71}. Trace normalization
is performed separately in each row, so \eqref{eq:cqchoi71} is exactly
trace preserving. The inherited zero-pivot proof, rational squared
coordinates and integer-square-root ratio calculation apply row by row.
Their finite masks and anchor headers have fixed cardinality; the row
body has at most $v_x$ signed grid coordinates. Multiplying these
counts gives the asserted exponent $V$. Only valid decoded words count.
The block-rank and zero statements hold row by row and pass to the direct
sum in \eqref{eq:cqchoi71}. Row $x$ has normalization denominator
$T_x\le h_xB^2$. A common denominator is the least common multiple of
the $T_x$, bounded by their product. This proves the final assertion.
\end{proof}

\begin{proof}[Proof of Theorem~\ref{thm:cqentropy71}]
For the upper bound take
$B=\max\{1,\lceil4\sqrt{NV}/\delta\rceil\}$. The last lemma gives
error at most $\delta$ and at most $C(2B+1)^V$ centres, hence the
upper description law for $\delta\le\delta_*<2$. Real targets admit
the same factor rounding construction; the decoded matrices are still
rational. For supplied rational targets, no irrational factor is stored:
the exact inherited square-comparison algorithm determines the digits.
Thus the real-family rational-centre upper bound does not restrict the
targets to a countable subclass.

For the lower bound choose, in every row with $v_x>0$, the gauge-fixed
maximal-rank factor patch used in Theorem~\ref{thm:rankentropy68}.
Its inverse is the local leading-block Cholesky map. Rows with $v_x=0$
are fixed. A Cartesian product of sufficiently small closed coordinate
balls contains a $V$-dimensional ball. The map to the tuple
$(\sigma_{xy})$, and hence to $J$, is bi-Lipschitz; direct sum is an
isometry for tuple Frobenius norm. This supplies one fixed
$V$-dimensional patch with no claim of a global factor chart.
For $\delta\le1/32$ pack it at separation greater than
$32d\delta/\sqrt N$. Corollary~\ref{cor:localchoi68} gives distances
larger than $2\delta$, so even an unrestricted legal centre covers at
most one packing point. This gives
$c(\sqrt N/\delta)^V$ centres. For $1/32<\delta\le\delta_*$ use
Theorem~\ref{thm:fullpacking71} on the same patch and the boundedness
of $\log(1/\delta)$ on that interval. These lower bounds apply also
to rational centres, which form a smaller class. Taking logarithms,
with the usual ceiling for fixed-length words, completes the proof.
If $V=0$, each nonempty row consists of a specified one-dimensional
unit block; its location is fixed by the public rank bounds. Thus the
whole family is a known singleton.
\end{proof}
For example, with scalar output $n=1$ and all rank bounds equal to
one, \eqref{eq:cqentropy71} has $V=d(m-1)$: it is a boundary coding
law for the rows of a classical stochastic matrix. With one quantum
output and full row ranks, $V=d(n^2-1)$; with one pure output per row,
$V=d(2n-2)$. Mixed and pure conditional states need not commute.
These dimensions are standard rank-stratum dimensions; the theorem
identifies their joint coding scale under the declared adaptive metric.
The measured input basis is fixed. No law for varying measurement bases,
unknown instruments, or all irreversible quantum dynamics follows from
this result.
